\section{November Musings}

There are those who think they can hide from the pains, frustrations, and sorrows of life simply by thinking positive thoughts. This is not the mark of high spiritual evolution, as they think, but rather the sign of a lack of depth.

There is the new age notion that our thoughts can entice the universe to satisfy all our desires, just as a baby commands the teat by crying. But even the baby needs to be weaned from milk and start eating flesh.

It is hard to believe that the end of desire will bring the end of suffering, when all our instincts convince us that suffering ends when desires are satisfied.

Most of those attracted to a spiritual life, in their initial excitement, try to adopt the dress, mannerisms, and speech of what they imagine a ``spiritual" person to be like. This is most certainly wrong. At first, change nothing, but begin to bring order to your own mind. Your life will then follow naturally.

I'm amazed at how many people show no sign of change after adopting some spiritual practice, even for a long time. Their views and ideas are indistinguishable from their quite ordinary neighbors.

Methodism was a decisive movement in the anti-tradition when John Wesley discerned God's Will by the ``warm feeling in the heart". For Tradition, such discernment is a matter of the Intelligence or Intellect, not ephemeral and unreliable feelings. This attitude dominates American culture where nearly everyone is convinced his opinion is correct because it is accompanied by a comforting or strong feeling.

This attitude is worsened by the idea of progress. The past is considered to have been inhabited by ignorant, superstitious, stupid, intolerant, and evil men. The corollary is that men of today have ``evolved", apparently through no effort of their own. Hence, they cannot possibly be ignorant, superstitious, stupid, intolerant, or evil, no matter what they say, do, or think.



\flrightit{Posted on 2012-11-12 by Cologero }

\begin{center}* * *\end{center}

\begin{footnotesize}\begin{sffamily}



\texttt{Michael on 2012-11-13 at 18:45 said: }

This reminds me of a question I've had for a while but I couldn't find the right place to ask it. Is there anything within esoteric tradition that addresses maintaining or reclaiming physical health?


\hfill

\texttt{Dominion on 2012-11-13 at 19:36 said: }

I can't think of any particular text, but it would seem that to maintain a healthy inner order would require its manifestation in a healthy outer form. The beauty of the outer form, in turn, inspires one to return to the source of the truths it reflects. The reflection of metaphysical principles is fairly obvious.


\hfill

\texttt{Anna Swadling on 2016-11-13 at 12:57 said: }

Michael dear, I personally witness a stronger presence of my spirit when my physical form is what I call a diamond body–when the two are one. I think I must have also read this from Evola some years ago. More thoughts welcome . . .


\end{sffamily}\end{footnotesize}
